
\section{\texorpdfstring{\Cref{chap:imitation_undelayed}}{} 的补充结果}


\subsection{涉及 Wasserstein 距离的界}


\begin{prop}
\label{pp:expected_wass_dist}
    考虑两个实值随机变量 $X,Y$，其分布分别为 $\eta,\nu$。则有：
    \begin{align*}
        \left\vert \expectedvalue[X]-\expectedvalue[Y]  \right\vert\leq \wassersteindistance(\eta\Vert \nu).
    \end{align*}
\end{prop}
\begin{proof}
    有，
    \begin{align}
        \expectedvalue[X]-\expectedvalue[Y]
        &= \int_{\realnumbers} x(\eta(x)-\nu(x))~dx
        \nonumber\\
        &\leq\sup_{\left\Vert f\right\Vert_L\leq 1}\left\vert\int_{\realnumbers} f(x)(\eta(x)-\nu(x))~dx\right\vert
        \label{eq:x_lip}\\
        &\leq \wassersteindistance(\eta\Vert \nu),\label{eq:recognise_wass}
    \end{align}
    其中 \Cref{eq:x_lip} 基于恒等函数是 1-LC 的事实，\Cref{eq:recognise_wass} 由识别 Wasserstein 距离得到。
    对 $\expectedvalue[Y]-\expectedvalue[X]$ 可进行相同的推理，Wasserstein 距离的对称性即完成证明。
\end{proof}

下一个结果断言，若将 $L$-LC 函数应用于两个随机变量，所得两个随机变量的分布之间的 Wasserstein 距离受原始 Wasserstein 距离乘以因子 $L$ 的约束。

\begin{prop}
\label{pp:wass_dist_lip_g}
    考虑度量空间 $\Omega$ 上的两个概率测度 $\eta$ 和 $\nu$，以及 $L_f$-LC 函数 $f:\Omega\rightarrow \realnumbers$。
    对于服从分布 $\eta$ 的某个随机变量 $X$，记 $f_\eta$ 为 $f(X)$ 的分布。类似定义 $f_\nu$。
    则有，
    \begin{align*}
        \wassersteindistance(f_{\eta}\Vert f_{\nu})\leq L_f \wassersteindistance(\eta \Vert \nu).
    \end{align*}
\end{prop}
\begin{proof}
    直接证明结果：
    \begin{align}
        \wassersteindistance(f_{\eta}\Vert f_{\nu})
        &= \sup_{\left\Vert g\right\Vert_L\leq 1}\left\vert\int_{\mathbb R} g(x)(f_{\eta}(x)-f_{\nu}(x)) ~dx\right\vert
        \label{eq:def_wass_f}
        \\
        &=\sup_{\left\Vert g\right\Vert_L\leq 1}
        \left\vert\int_{\mathbb R} g(x)f_{\eta}(x)~dx-\int_{\mathbb R} g(x)f_{\nu}(x) dx\right\vert,\nonumber
        \\
        &= \sup_{\left\Vert g\right\Vert_L\leq 1}
        \left\vert\int_{\mathbb R} g(f(x))\eta(x) ~dx-\int_{\mathbb R} g(f(x)) \nu(x) dx\right\vert\label{eq:def_eta_nu}
        \\
        &= \sup_{\left\Vert g\right\Vert_L\leq 1}
        \left\vert \int_{\mathbb R} g(f(x))(\eta(x) - \nu(x)) ~dx\right\vert,\nonumber
    \end{align}
    其中 \Cref{eq:def_wass_f} 由 Wasserstein 距离的定义成立，\Cref{eq:def_eta_nu} 由 $f_\eta$ 和 $f_\nu$ 的定义成立。结论由观察到 $g(f(x))$ 是 $L_f$-LC 函数（因为 1-LC 函数与 $L_f$-LC 函数的复合）得到。
\end{proof}

\begin{prop}
\label{pp:expected_q_bound}
    设 $\markovdecisionprocess$ 为一个 \gls{mdp}，$\pi$ 为一个策略，其 Q 函数在第二个参数上是 $L_{Q}$-LC 的。
    则对于 $\actionspace$ 上的任意两个概率分布 $\eta,\nu$，
    \begin{align*}
        \left\vert \expectedvalue_{\substack{X\sim \eta\\Y\sim \nu}}[Q^{\pi}(s,X)-Q^{\pi}(s,Y)] \right\vert
        \leq L_{Q} \wassersteindistance(\eta(\cdot)\Vert\nu(\cdot)).
    \end{align*}
\end{prop}
\begin{proof}
    对于某个 $s\in\statespace$，设 $f_{\eta}(X)$ 和 $f_{\nu}(X)$ 分别为 $Q^{\pi}(s,X)$ 和 $Q^{\pi}(s,Y)$ 的分布。
    首先，应用 \Cref{pp:expected_wass_dist} 得到，
    \begin{align*}
        \left\vert \expectedvalue_{\substack{X\sim \eta\\Y\sim \nu}}[Q^{\pi}(s,X)-Q^{\pi}(s,Y)] \right\vert
        \leq  \wassersteindistance(g_{\eta}\Vert g_{\nu}).
    \end{align*}
    然后，应用 \Cref{pp:wass_dist_lip_g} 得到，
    \begin{align*}
        \wassersteindistance(f_{\eta}\Vert f_{\nu})
        \leq L_{Q} \wassersteindistance(\eta\Vert \nu).
    \end{align*}
\end{proof}


\begin{prop}
\label{pp:lip_dmdp}
    设 $\markovdecisionprocess$ 是一个 $(L_P,L_r)$-LC \gls{mdp}。
    设 $\delayedmarkovdecisionprocess$ 为基于 $\markovdecisionprocess$ 构建的常值 $\delay$-延迟 \gls{dmdp}，其转移函数为 $\augmentedtransitionfunction$，奖励函数为 $\augmentedexpectedrewardfunction$。
    则 $\delayedmarkovdecisionprocess$ 是 $(\max(1,L_P),L_r)$-LC 的。
\end{prop}
\begin{proof}
    首先分析奖励函数。
    设 $x=(s,a_1,\dots,a_\delay)$ 和 $x'=(s','a_1,\dots,a_\delay')$ 属于 $\augmentedstatespace$，$a,a'\in\actionspace$，有，
    \begin{align*}
        \lvert \augmentedexpectedrewardfunction(x,a)-\augmentedexpectedrewardfunction(x',a')\rvert
        &= \lvert \expectedrewardfunction(s,a_1)-\expectedrewardfunction(s',a_1')\rvert
        \\
        &\le
        L_r \left( \distance_{\statespace}(s,s') + \distance_{\actionspace}(a_1,a_1')\right)
        \\
        &\le
        L_r \left( \distance_{\augmentedstatespace}(x,x') + \distance_{\actionspace}(a,a')\right),
    \end{align*}
    其中使用了 \Cref{eq:delayed_reward_deterministic}。
    对于转移函数，为方便起见记 $a_{\delay+1}=a$ 和 $a_{\delay+1}'=a'$。有，
    \begin{align*}
        \wassersteindistance[1]&(\augmentedtransitionfunction(\cdot\vert x,a_{\delay+1})\Vert \augmentedtransitionfunction(\cdot\vert x',a_{\delay+1}'))
        \\
        &= \sup_{\left\Vert f\right\Vert_L\leq 1} \left\vert \int_{\augmentedstatespace} f(x'')(\augmentedtransitionfunction(x''\vert x,a_{\delay+1})-\augmentedtransitionfunction(x''\vert x',a_{\delay+1}'))\; (dx'') \right\vert
        \\
        &= \sup_{\left\Vert f\right\Vert_L\leq 1} \left\vert \int_{\augmentedstatespace} f(x'')\left[\transitionfunction(s''\vert s,a_{1})\prod_{i=1}^{\delay}\delta_{a_{i+1}}(a_i'')-\transitionfunction(s''\vert s',a_{1}')\prod_{i=1}^{\delay}\delta_{a_{i+1}}(a_i'')\right]\; (dx'') \right\vert
        \\
        &\le \sup_{\left\Vert f\right\Vert_L\leq 1} \left\vert \underbrace{\int_{\augmentedstatespace} f(x'')\prod_{i=1}^{\delay}\delta_{a_{i+1}}(a_i'')\left[\transitionfunction(s''\vert s,a_{1})-\transitionfunction(s''\vert s',a_{1}')\right]\; (dx'')}_{A} \right\vert
        \\
        &\qquad + \sup_{\left\Vert f\right\Vert_L\leq 1} \left\vert \underbrace{\int_{\augmentedstatespace} f(x'')\transitionfunction(s''\vert s',a_{1}') \left[\prod_{i=1}^{\delay}\delta_{a_{i+1}}(a_i'')-\prod_{i=1}^{\delay}\delta_{a_{i+1}}(a_i'')\right]\; (dx'')}_{B} \right\vert,
    \end{align*}
    其中在最后一个不等式中加减了相同量，并使用记号 $x''=(s'',a_1'',\dots,a_\delay'')$。
    考虑第一项，
    \begin{align*}
        A &= \int_{\actionspace^\delay} \prod_{i=1}^{\delay}\delta_{a_{i+1}}(a_i'') \int_{\statespace}f(x'')\left[\transitionfunction(s''\vert s,a_{1})-\transitionfunction(s''\vert s',a_{1}')\right]\;(ds'' da_1'' \dots da_{\delay}'')
        \\
        &\le L_P(\distance_{\statespace}(s,s') + \distance_{\actionspace}(a_1,a_1')) \int_{\actionspace^{\delay}} \prod_{i=1}^{\delay}\delta_{a_{i+1}}(a_i'')  \;(da_1'' \dots da_{\delay}'')
        \\
        &= L_P(\distance_{\statespace}(s,s') + \distance_{\actionspace}(a_1,a_1')),
    \end{align*}
    其中利用了 $f$ 在 $\augmentedstatespace$ 上是 1-LC 的，因此 $f$ 在 $\statespace$ 上也是 1-LC 的。
    对于另一项，类似地，
    \begin{align*}
        B &=\int_{\statespace}\transitionfunction(s''\vert s',a_{1}') \int_{\actionspace^\delay}f(x'') \left[\prod_{i=1}^{\delay}\delta_{a_{i+1}}(a_i'')-\prod_{i=1}^{\delay}\delta_{a_{i+1}}(a_i'')\right]\;\;(da_1'' \dots da_{\delay}''ds'')
        \\
        &\leq \sum_{i=2}^{\delay+1}\distance_{\actionspace}(a_{i},a_{i}').
    \end{align*}
    合并各项，有，
    \begin{align*}
        \wassersteindistance[1](\augmentedtransitionfunction(\cdot\vert x,a_{\delay+1})\Vert \augmentedtransitionfunction(\cdot\vert x',a_{\delay+1}'))
        &\le L_P(\distance_{\statespace}(s,s') + \distance_{\actionspace}(a_1,a_1')) + \sum_{i=2}^{\delay+1}\distance_{\actionspace}(a_{i},a_{i}')
        \\
        &\le \max(1,L_P) (\distance_{\statespace}(x,x') + \distance_{\actionspace}(a_{\delay+1},a_{\delay+1}'))
    \end{align*}
\end{proof}


\begin{prop}
\label{pp:lip_dida_pol}
    设 $\markovdecisionprocess$ 是一个 $(L_T)$-TLC \gls{mdp}。
    设 $\delayedmarkovdecisionprocess$ 为基于 $\markovdecisionprocess$ 构建的常值 $\delay$-延迟 \gls{dmdp}。
    假设 $\delayedmarkovdecisionprocess$ 的策略 $\delayedpolicy$ 满足 \Cref{eq:belief_pol}，其中 $\pi$ 是 $\markovdecisionprocess$ 中的 $L_\pi$-LC 策略。
    则 $\delayedpolicy$ 是 $2\delay L_\pi L_T$-LC 的。
\end{prop}
\begin{proof}
    设 $x=(s,a_1,\dots,a_\delay)$ 和 $x'=(s','a_1,\dots,a_\delay')$ 属于 $\augmentedstatespace$。
    \begin{align}
        \wassersteindistance[1](\delayedpolicy(\cdot\vert x)\Vert \delayedpolicy(\cdot\vert x'))
        &= \sup_{\left\Vert f\right\Vert_L\leq 1} \left\vert \int_{\actionspace} f(a)\pi(a\vert s)(\augmentedbelief(s\vert x)-\augmentedbelief(s\vert x'))\; (da) \right\vert
        \nonumber\\
        &\le L_\pi\sup_{\left\Vert f\right\Vert_L\leq 1} \left\vert \int_{\statespace} f(s)(\augmentedbelief(s\vert x)-\augmentedbelief(s\vert x'))\; (da) \right\vert
        \label{eq:lpi_in_s}
        \\
        &\le 2 L_\pi L_T,
        \label{eq:lpilc}
    \end{align}
    其中在 \Cref{eq:lpi_in_s} 中使用了 $s\mapsto f(a)\pi(a\vert s)$ 是 $L_\pi$-LC 的，在 \Cref{eq:lpilc} 中使用了 \Cref{lem:bound_sigma_tlc}。
\end{proof}

\subsection{对 \texorpdfstring{$\sigma_b^{\rho}$}{}} 的界定
在本小节中，提供两种方式来界定来自 \Cref{subsec:imitation_upper_bound} 的项 $\textstyle\sigma_b^\rho = \expectedvalue_{\substack{x'\sim \delayeddiscountedstateoccupancydistribution[x][\delayedpolicy]\\ s,s'\sim \augmentedbelief(\cdot\vert x')}}\left[ \distance_{\statespace}(s,s') \right]$，其中 $\rho$ 是 $\augmentedstatespace$ 上的分布。
对于第一种方式，假设 $\statespace\subset\realnumbers^n$ 配备了欧几里得范数。

\begin{lemma}[欧几里得界]
\label{lem:bound_sigma_eucl}
    设 $\markovdecisionprocess$ 为一个 \gls{mdp}，其中 $\statespace\subset\realnumbers^n$ 配备了欧几里得范数。则有，
    \begin{align*}
        \sigma_b^\rho \leq \sqrt{2}\expectedvalue_{x'\sim \delayeddiscountedstateoccupancydistribution[\rho][\delayedpolicy](\cdot)}\left[\sqrt{ \variance_{s\sim \augmentedbelief(\cdot|x')}(s|x')}\right],
    \end{align*}
    其中 $\rho$ 是 $\augmentedstatespace$ 上的分布。
\end{lemma}
\begin{proof}
    证明基于以下步骤，
    \begin{align}
        \sigma_b^\rho
        &= \expectedvalue_{\substack{x'\sim \delayeddiscountedstateoccupancydistribution[\rho][\delayedpolicy]\\ s,s'\sim \augmentedbelief(\cdot\vert x')}}\left[ \distance_{\statespace}(s,s') \right]\nonumber
        \\
        &= \expectedvalue_{\substack{x'\sim \delayeddiscountedstateoccupancydistribution[\rho][\delayedpolicy]\\ s,s'\sim \augmentedbelief(\cdot\vert x')}}\left[ \sqrt{(s'-s)^2} \right]\label{eq:def_euclid}
        \\
        &= \expectedvalue_{x'\sim \delayeddiscountedstateoccupancydistribution[\rho][\delayedpolicy]}\sqrt{\expectedvalue_{s,s'\sim \augmentedbelief(\cdot\vert x')} \left[ (s'-s)^2 \right]}\label{eq:jensen_var}
        \\
        &= \expectedvalue_{x'\sim \delayeddiscountedstateoccupancydistribution[\rho][\delayedpolicy]}\left[\sqrt{ \variance_{s\sim \augmentedbelief(\cdot|x')}(s|x')}\right].\label{eq:idd_var_s},
    \end{align}
    其中 \Cref{eq:def_euclid} 显式写出了欧几里得范数，\Cref{eq:jensen_var} 应用了 Jensen 不等式，
    \Cref{eq:idd_var_s} 基于如下事实：若 $s'$ 和 $s$ 是独立同分布的，则 $\textstyle\expectedvalue_{s,s'\sim \augmentedbelief(\cdot\vert x')} \left[ (s'-s)^2 \right] = 2\variance_{s\sim \augmentedbelief(\cdot|x')}[s]$。
\end{proof}


现在提供第二个结果，它假设 \gls{mdp} 是时间-Lipschitz 的（见 \Cref{def:time_lip} 和 \Cref{eq:time_lip_non_int}）。
在提供结果之前，先证明以下中间结果。

\begin{prop}
\label{pp:tlc_delay}
    设 $\markovdecisionprocess$ 为一个 $L_T$-TLC \gls{mdp}\footnote{注意，由于考虑非整数延迟，我们在 \Cref{eq:time_lip_non_int} 中扩展了 TLC 假设。}。设 $x=(s_1,a_1,\dots,a_\delay)\in\statespace\times\actionspace^\delay$ 为延迟 $\delay\in\realnumbers[\ge0]$ 下的增广状态。
    则有：
    \begin{align*}
        \wassersteindistance\left(\augmentedbelief(\cdot\vert x)\Vert \delta_{s_1} \right)\leq \delay L_T
    \end{align*}
\end{prop}
\begin{proof}
    首先证明整数延迟的结果，再证明非整数延迟的结果。

    \noindent\textbf{整数延迟。}\indent 设 $\delay\in\naturalnumbers$。
    采用归纳法证明。
    初始化时，当 $\delay=0$，结果显然成立，因为当前状态对智能体是已知的。
    注意 $\delay=1$ 由 $L_T$-TLC 假设成立。
    对于递推步骤，假设该断言对某个 $\delay \in\naturalnumbers$ 成立。设 $x=(s_1,a_1,\dots.a_{\delay+1})\in\statespace\times\actionspace^{\delay+1}$。则，
    \begin{align}
        &\wassersteindistance\left(\augmentedbelief[\delay+1](\cdot\vert x)\Vert \delta_{s_1} \right)
        \\
        &= \sup_{\left\Vert f\right\Vert_L\leq 1}\left\vert\int_{\statespace} f(s')\left( \augmentedbelief[\delay+1](s' \vert x) -\delta_{s_1}(s')\right)ds' \right\vert\nonumber
        \\
        &= \sup_{\left\Vert f\right\Vert_L\leq 1}\left\vert\int_{\statespace} p(s_2\vert s_1,a_1) \int_{\statespace} f(s')\left( \augmentedbelief(s' \vert s_2,a_2,\cdots,a_\delay) -\delta_{s_1}(s')\right)ds' \right\vert
        \label{eq:cond_s2}
        \\
        &= \sup_{\left\Vert f\right\Vert_L\leq 1}\left\vert\int_{\statespace} p(s_2\vert s_1,a_1) \int_{\statespace} f(s')\left( \augmentedbelief(s' \vert s_2,a_2,\cdots,a_\delay) -\delta_{s_2}(s')\right)ds' \right.\nonumber
        \\
        &\qquad + \left.\int_{\statespace} p(s_2\vert s_1,a_1) \int_{\statespace} f(s')\left( \delta_{s_2}(s) -\delta_{s_1}(s')\right)ds' \right\vert
        \label{eq:add_remove_delta}
        \\
        &\leq \underbrace{\sup_{\left\Vert f\right\Vert_L\leq 1}\left\vert\int_{\statespace} p(s_2\vert s_1,a_1) \int_{\statespace} f(s')\left( \augmentedbelief(s' \vert s_2,a_2,\cdots,a_\delay) -\delta_{s_2}(s')\right)ds'  \right\vert}_A\nonumber
        \\
        &\qquad + \underbrace{\wassersteindistance\left(P(\vert s_1,a_1)\Vert \delta_{s_1}\right)}_B
        \nonumber,
    \end{align}
    其中 \eqref{eq:cond_s2} 通过对下一个观测状态 $s_2$ 取条件得到，\Cref{eq:add_remove_delta} 通过加减相同量 $\delta_{s_2}$ 得到。

    项 $A$ 可使用 $\delay$ 处的断言进行界定，而 $B$ 直接由 TLC 假设得到。
    因此有，
    \begin{align*}
        \wassersteindistance\left(\augmentedbelief[\delay+1](\cdot\vert x)\Vert \delta_{s_1} \right) \leq (\delay+1) L_T.
    \end{align*}
    通过归纳法，证明了该结果对任意 $\delay\in\naturalnumbers$ 成立。

    \noindent\textbf{非整数延迟。}\indent 为证明结果在更一般的 $\delay\in\realnumbers[\ge0]$ 情况下成立，可以将延迟分解为小数部分（$\fractionalpart{\delay}$）和整数部分（$\integerpart{\delay}$）。应用与之前类似的推理，会得到类似的项 $A$ 和 $B$，其中 $A$ 涉及延迟的整数部分（该断言对其成立），$B$ 涉及 $\delay$ 的小数部分（该断言在 \Cref{eq:time_lip_non_int} 下对其成立）。
\end{proof}



\begin{lemma}[时间-Lipschitz 界]
\label{lem:bound_sigma_tlc}
    设 $\markovdecisionprocess$ 为一个 $L_T$-TLC \gls{mdp}。考虑从 $\markovdecisionprocess$ 获得的 $\delay$-延迟 \gls{dmdp} $\delayedmarkovdecisionprocess$，其中 $\delay\in\realnumbers[\ge0]$。则，
    \begin{align*}
        \sigma_b^\rho  \leq 2\delay L_T,
    \end{align*}
    其中 $\rho$ 是 $\augmentedstatespace$ 上的分布。
\end{lemma}
\begin{proof}
    对于增广状态 $x$，设 $s_x$ 为其包含的最近观测状态。
    则可以写出以下内容，
    \begin{align}
        \sigma_b^\rho
        &= \expectedvalue_{\substack{x\sim \delayeddiscountedstateoccupancydistribution[\rho][\delayedpolicy]\\ s,s'\sim \augmentedbelief(\cdot\vert x)}}\left[ \distance_{\statespace}(s,s') \right]\nonumber
        \\
        &\leq \expectedvalue_{\substack{x\sim \delayeddiscountedstateoccupancydistribution[\rho][\delayedpolicy]\\ s\sim \augmentedbelief(\cdot\vert x')}}\left[ \distance_{\statespace}(s,s_x) \right] + \expectedvalue_{\substack{x\sim \delayeddiscountedstateoccupancydistribution[\rho][\delayedpolicy]\\ s'\sim \augmentedbelief(\cdot\vert x)}}\left[ \distance_{\statespace}(s_x,s') \right]\label{eq:triang_ineq_tlc}
        \\
        &= 2\expectedvalue_{x\sim \delayeddiscountedstateoccupancydistribution[\rho][\delayedpolicy]}\left[\int_{\statespace} d_{\statespace}(s,s_x) \augmentedbelief(s\vert x)~ds\right] \nonumber
        \\
        &= 2\expectedvalue_{x\sim \delayeddiscountedstateoccupancydistribution[\rho][\delayedpolicy]}\left[\int_{\statespace} \distance_{\statespace}(s,s_x) \left(\augmentedbelief(s\vert x) - \delta_{s_x}(s) \right)~ds\right]\label{eq:null_term}
        \\
        &\leq 2 \expectedvalue_{x\sim \delayeddiscountedstateoccupancydistribution[\rho][\delayedpolicy]}\left[\wassersteindistance\left(\augmentedbelief(\cdot\vert x)\Vert \delta_{s_x} \right)\right]\label{eq:wass_recognise},
    \end{align}
    其中 \Cref{eq:triang_ineq_tlc} 由三角不等式成立；在 \Cref{eq:null_term} 中添加了 $\textstyle\int_{\statespace}d_{\statespace}(s,s_x)\delta_{s_x}(s)ds=0$ 项；在 \Cref{eq:wass_recognise} 中使用了 Wasserstein 距离的定义。
    为完成证明，在期望内应用 \Cref{pp:tlc_delay}。
\end{proof}